%%%%%%%%%%%%%%%%%%%%%%%%%%%%%%%%%%%%%%%%%%%%%%%%%%%%%%%%%%%%%%%%%%%%%%%%%%%
%% Hydra CS Lab -- Development Environment Setup Guide
%% SUNY New Paltz Computer Science Department
%%%%%%%%%%%%%%%%%%%%%%%%%%%%%%%%%%%%%%%%%%%%%%%%%%%%%%%%%%%%%%%%%%%%%%%%%%%
\documentclass[11pt,a4paper]{article}

% ---------- packages ----------
\usepackage[margin=1in]{geometry}
\usepackage[utf8]{inputenc}
\usepackage[T1]{fontenc}
\usepackage{mathpazo}            % Palatino serif (body + math)
\usepackage[scaled=0.85]{inconsolata}  % clean monospace for code
\usepackage{microtype}           % better character spacing/kerning
\usepackage{graphicx}
\usepackage{xcolor}
\usepackage{listings}
\usepackage{enumitem}
\usepackage{booktabs}
\usepackage{longtable}
\usepackage{array}
\usepackage{tabularx}
\usepackage{fancyhdr}
\usepackage{titlesec}
\usepackage[colorlinks=true,
            linkcolor=blue!70!black,
            urlcolor=blue!60!black,
            citecolor=blue!70!black,
            bookmarks=true,
            pdfauthor={SUNY New Paltz CS Department},
            pdftitle={Hydra CS Lab -- Development Environment Setup Guide}]
            {hyperref}

% ---------- colours ----------
\definecolor{codebg}{RGB}{245,245,245}
\definecolor{codeframe}{RGB}{200,200,200}
\definecolor{codekw}{RGB}{0,0,180}
\definecolor{codecomment}{RGB}{0,128,0}
\definecolor{codestr}{RGB}{163,21,21}
\definecolor{sectionblue}{RGB}{0,51,102}
\definecolor{notebg}{RGB}{255,249,230}
\definecolor{noteframe}{RGB}{220,190,100}
\definecolor{warnbg}{RGB}{255,240,240}
\definecolor{warnframe}{RGB}{200,80,80}

% ---------- listings ----------
\lstset{
    backgroundcolor=\color{codebg},
    frame=single,
    rulecolor=\color{codeframe},
    basicstyle=\ttfamily\small,
    keywordstyle=\color{codekw}\bfseries,
    commentstyle=\color{codecomment}\itshape,
    stringstyle=\color{codestr},
    breaklines=true,
    breakatwhitespace=false,
    showstringspaces=false,
    tabsize=2,
    captionpos=b,
    xleftmargin=1em,
    xrightmargin=1em,
    aboveskip=1em,
    belowskip=0.8em,
    literate={~}{{\raisebox{0.5ex}{\texttildelow}}}{1},
}

\lstdefinestyle{bash}{
    language=bash,
    morekeywords={sudo,apt,brew,choco,curl,docker,kubectl,k3d,make,git,
                  cd,mkdir,cat,echo,chmod,wget,ssh,export},
}

\lstdefinestyle{yaml}{
    basicstyle=\ttfamily\small,
    morecomment=[l]{\#},
    morestring=[b]',
    morestring=[b]",
}

% ---------- custom commands ----------
\newcommand{\notebox}[1]{%
    \vspace{0.5em}
    \noindent
    \fcolorbox{noteframe}{notebg}{%
        \begin{minipage}{0.94\textwidth}
            \textbf{Note:} #1
        \end{minipage}%
    }
    \vspace{0.5em}
}

\newcommand{\warnbox}[1]{%
    \vspace{0.5em}
    \noindent
    \fcolorbox{warnframe}{warnbg}{%
        \begin{minipage}{0.94\textwidth}
            \textbf{Warning:} #1
        \end{minipage}%
    }
    \vspace{0.5em}
}

\newcommand{\cmdref}[1]{\texttt{#1}}

% ---------- header / footer ----------
\pagestyle{fancy}
\fancyhf{}
\fancyhead[L]{\small Hydra CS Lab -- Dev Environment Guide}
\fancyhead[R]{\small SUNY New Paltz}
\fancyfoot[C]{\thepage}
\renewcommand{\headrulewidth}{0.4pt}
\renewcommand{\footrulewidth}{0.2pt}

% ---------- section styling ----------
\titleformat{\section}
  {\normalfont\Large\bfseries\color{sectionblue}}
  {\thesection}{1em}{}
\titleformat{\subsection}
  {\normalfont\large\bfseries\color{sectionblue!80!black}}
  {\thesubsection}{1em}{}
\titleformat{\subsubsection}
  {\normalfont\normalsize\bfseries\color{sectionblue!60!black}}
  {\thesubsubsection}{1em}{}

%%%%%%%%%%%%%%%%%%%%%%%%%%%%%%%%%%%%%%%%%%%%%%%%%%%%%%%%%%%%%%%%%%%%%%%%%%%
%% TITLE PAGE
%%%%%%%%%%%%%%%%%%%%%%%%%%%%%%%%%%%%%%%%%%%%%%%%%%%%%%%%%%%%%%%%%%%%%%%%%%%
\begin{document}

\begin{titlepage}
    \centering
    \vspace*{3cm}

    {\Huge\bfseries\color{sectionblue}
        Hydra CS Lab\\[0.3em]
        Development Environment\\[0.3em]
        Setup Guide
    }

    \vspace{1.5cm}

    {\Large\color{gray}
        SUNY New Paltz\\
        Computer Science Department
    }

    \vspace{2cm}

    \rule{0.6\textwidth}{0.5pt}

    \vspace{1cm}

    {\large
        \textbf{Production Cluster}\\[0.5em]
        \begin{tabular}{l l l}
            Hydra & 192.168.1.160 & Control Plane \\
            Chimera & 192.168.1.150 & GPU Inference \\
            Cerberus & 192.168.1.242 & GPU Training \\
        \end{tabular}
    }

    \vspace{2cm}

    {\large \today}

    \vfill

    {\small
        Repository: \url{https://github.com/compsci-suny-newpaltz/hydra-saml-auth}
    }
\end{titlepage}

%%%%%%%%%%%%%%%%%%%%%%%%%%%%%%%%%%%%%%%%%%%%%%%%%%%%%%%%%%%%%%%%%%%%%%%%%%%
%% TABLE OF CONTENTS
%%%%%%%%%%%%%%%%%%%%%%%%%%%%%%%%%%%%%%%%%%%%%%%%%%%%%%%%%%%%%%%%%%%%%%%%%%%
\newpage
\tableofcontents
\newpage

%%%%%%%%%%%%%%%%%%%%%%%%%%%%%%%%%%%%%%%%%%%%%%%%%%%%%%%%%%%%%%%%%%%%%%%%%%%
%% 1. PREREQUISITES
%%%%%%%%%%%%%%%%%%%%%%%%%%%%%%%%%%%%%%%%%%%%%%%%%%%%%%%%%%%%%%%%%%%%%%%%%%%
\section{Prerequisites}
\label{sec:prerequisites}

Before setting up the Hydra development environment, you need the following
tools installed on your machine. Instructions are provided for each supported
platform.

\subsection{Required Software}

\begin{itemize}[leftmargin=2em]
    \item \textbf{Docker} -- container runtime (Docker Desktop or Docker Engine)
    \item \textbf{kubectl} -- Kubernetes CLI for cluster interaction
    \item \textbf{k3d} -- lightweight Kubernetes wrapper around k3s in Docker
    \item \textbf{Git} -- version control
    \item \textbf{Make} -- build automation (GNU Make)
    \item \textbf{Node.js 18+} -- for local development outside containers (optional)
\end{itemize}

%% ---- Windows ----
\subsection{Windows}
\label{sec:prereq-windows}

\subsubsection{Install WSL2}

Windows Subsystem for Linux 2 is required for Docker and the development
workflow. Open PowerShell as Administrator:

\begin{lstlisting}[style=bash]
# Install WSL2 with Ubuntu 22.04
wsl --install -d Ubuntu-22.04

# Verify installation
wsl -l -v
\end{lstlisting}

After installation, restart your computer and complete the Ubuntu setup
(username and password) when prompted.

\subsubsection{Install Docker Desktop}

\begin{enumerate}[leftmargin=2em]
    \item Download Docker Desktop from
          \url{https://www.docker.com/products/docker-desktop/}
    \item During installation, ensure \textbf{``Use WSL 2 based engine''} is
          checked.
    \item After installation, open Docker Desktop and go to
          \textbf{Settings $\rightarrow$ Resources $\rightarrow$ WSL Integration}.
    \item Enable integration with your Ubuntu-22.04 distribution.
    \item Click \textbf{Apply \& Restart}.
\end{enumerate}

Verify from a WSL2 Ubuntu terminal:

\begin{lstlisting}[style=bash]
docker --version
docker run hello-world
\end{lstlisting}

\subsubsection{Install kubectl, k3d, Git, and Make}

Using Chocolatey (run PowerShell as Administrator):

\begin{lstlisting}[style=bash]
# Install Chocolatey if not already present
Set-ExecutionPolicy Bypass -Scope Process -Force
[System.Net.ServicePointManager]::SecurityProtocol = `
    [System.Net.SecurityProtocolType]::Tls12
iex ((New-Object System.Net.WebClient).DownloadString(
    'https://community.chocolatey.org/install.ps1'))

# Install tools
choco install kubernetes-cli k3d git make -y
\end{lstlisting}

Alternatively, install inside WSL2 Ubuntu (see the Linux instructions in
Section~\ref{sec:prereq-linux}).

\subsubsection{Configure WSL2 Memory Limits}

Create or edit \cmdref{\%USERPROFILE\%\textbackslash.wslconfig} in Windows:

\begin{lstlisting}[style=yaml]
[wsl2]
memory=8GB
processors=4
swap=2GB
\end{lstlisting}

Restart WSL2: \cmdref{wsl --shutdown} then reopen your terminal.

\subsubsection{Clone the Repository}

From inside WSL2 Ubuntu:

\begin{lstlisting}[style=bash]
cd ~
git clone https://github.com/compsci-suny-newpaltz/hydra-saml-auth.git
cd hydra-saml-auth
\end{lstlisting}

%% ---- macOS ----
\subsection{macOS}
\label{sec:prereq-macos}

\subsubsection{Install Homebrew}

If Homebrew is not installed:

\begin{lstlisting}[style=bash]
/bin/bash -c "$(curl -fsSL \
    https://raw.githubusercontent.com/Homebrew/install/HEAD/install.sh)"
\end{lstlisting}

\subsubsection{Install Docker Desktop for Mac}

\begin{enumerate}[leftmargin=2em]
    \item Download Docker Desktop from
          \url{https://www.docker.com/products/docker-desktop/}
    \item Install the \texttt{.dmg} and drag Docker to Applications.
    \item Open Docker Desktop and allow it to finish initializing.
    \item Allocate at least 4 CPU cores and 8 GB RAM in
          \textbf{Settings $\rightarrow$ Resources}.
\end{enumerate}

\subsubsection{Install CLI Tools}

\begin{lstlisting}[style=bash]
brew install kubectl k3d git make
\end{lstlisting}

Verify installations:

\begin{lstlisting}[style=bash]
docker --version
kubectl version --client
k3d version
git --version
make --version
\end{lstlisting}

\subsubsection{Clone the Repository}

\begin{lstlisting}[style=bash]
cd ~
git clone https://github.com/compsci-suny-newpaltz/hydra-saml-auth.git
cd hydra-saml-auth
\end{lstlisting}

%% ---- Linux ----
\subsection{Linux (Ubuntu / Debian)}
\label{sec:prereq-linux}

\subsubsection{Install Docker}

\begin{lstlisting}[style=bash]
# Install Docker
sudo apt update
sudo apt install -y docker.io

# Add your user to the docker group (avoids needing sudo)
sudo usermod -aG docker $USER

# Apply group changes (or log out and back in)
newgrp docker

# Verify
docker run hello-world
\end{lstlisting}

\subsubsection{Install kubectl}

\begin{lstlisting}[style=bash]
# Download the latest stable kubectl
curl -LO "https://dl.k8s.io/release/$(curl -L -s \
    https://dl.k8s.io/release/stable.txt)/bin/linux/amd64/kubectl"
chmod +x kubectl
sudo mv kubectl /usr/local/bin/

# Verify
kubectl version --client
\end{lstlisting}

\subsubsection{Install k3d}

\begin{lstlisting}[style=bash]
curl -s https://raw.githubusercontent.com/k3d-io/k3d/main/install.sh \
    | bash

# Verify
k3d version
\end{lstlisting}

\subsubsection{Install Make and Git}

\begin{lstlisting}[style=bash]
sudo apt install -y make git
\end{lstlisting}

\subsubsection{Clone the Repository}

\begin{lstlisting}[style=bash]
cd ~
git clone https://github.com/compsci-suny-newpaltz/hydra-saml-auth.git
cd hydra-saml-auth
\end{lstlisting}

%%%%%%%%%%%%%%%%%%%%%%%%%%%%%%%%%%%%%%%%%%%%%%%%%%%%%%%%%%%%%%%%%%%%%%%%%%%
%% 2. QUICK START (DOCKER COMPOSE)
%%%%%%%%%%%%%%%%%%%%%%%%%%%%%%%%%%%%%%%%%%%%%%%%%%%%%%%%%%%%%%%%%%%%%%%%%%%
\section{Quick Start (Docker Compose)}
\label{sec:quickstart}

The fastest way to get the Hydra development environment running is via
Docker Compose. This mode uses \cmdref{dev/docker-compose.dev.yml} and
provides a mock SAML IdP, mock GPU metrics for Chimera and Cerberus, a
Traefik reverse proxy, OpenWebUI, and Ollama.

\subsection{Setup and Start}

\begin{lstlisting}[style=bash]
cd hydra-saml-auth/dev

# One-time setup: generates JWT keys, creates .env,
# configures /etc/hosts entries
make setup

# Start all services in the background
make start
\end{lstlisting}

\notebox{The \cmdref{make setup} target will attempt to add entries to your
hosts file (\cmdref{/etc/hosts} on Linux/macOS,
\cmdref{C:\textbackslash Windows\textbackslash System32\textbackslash
drivers\textbackslash etc\textbackslash hosts} on Windows). This requires
\cmdref{sudo} or Administrator privileges. If it fails, add the entries
manually (see Section~\ref{sec:troubleshooting}).}

\subsection{Accessible Services}

Once running, the following services are available:

\begin{center}
\begin{tabular}{l l}
    \toprule
    \textbf{Service} & \textbf{URL} \\
    \midrule
    Hydra Dashboard & \url{http://hydra.local/dashboard} \\
    OpenWebUI (Chat) & \url{http://gpt.hydra.local} \\
    Mock SAML IdP & \url{http://localhost:8080} \\
    Traefik Dashboard & \url{http://localhost:8081} \\
    \bottomrule
\end{tabular}
\end{center}

\subsection{Test Credentials}

The mock SAML IdP (SimpleSAMLphp) comes preconfigured with test users:

\begin{center}
\begin{tabular}{l l l}
    \toprule
    \textbf{Username} & \textbf{Password} & \textbf{Role} \\
    \midrule
    user1 & user1pass & Student \\
    \bottomrule
\end{tabular}
\end{center}

Log in at \url{http://hydra.local/dashboard} and use these credentials when
redirected to the mock IdP.

\subsection{Stopping the Environment}

\begin{lstlisting}[style=bash]
# Stop all services
make stop

# Stop and remove all volumes (clean slate)
make clean
\end{lstlisting}

%%%%%%%%%%%%%%%%%%%%%%%%%%%%%%%%%%%%%%%%%%%%%%%%%%%%%%%%%%%%%%%%%%%%%%%%%%%
%% 3. ARCHITECTURE OVERVIEW
%%%%%%%%%%%%%%%%%%%%%%%%%%%%%%%%%%%%%%%%%%%%%%%%%%%%%%%%%%%%%%%%%%%%%%%%%%%
\section{Architecture Overview}
\label{sec:architecture}

The development environment mirrors the three-node production cluster.
Below is the mapping between production components and their dev
equivalents in both Docker Compose mode and k3d mode.

\subsection{Node Mapping}

\begin{center}
\small
\begin{tabularx}{\textwidth}{l X X}
    \toprule
    \textbf{Production} & \textbf{Dev (Docker Compose)} & \textbf{Dev (k3d)} \\
    \midrule
    Hydra (control plane, 192.168.1.160) --- 256 GB RAM, 64 cores &
        hydra-saml-auth container &
        k3d-hydra-dev-server-0 \\
    \addlinespace
    Chimera (GPU inference, 192.168.1.150) --- 3x RTX 3090, 72 GB VRAM &
        mock-chimera + ollama + open-webui containers &
        k3d-hydra-dev-agent-0 (label: \texttt{gpu-role=inference}) \\
    \addlinespace
    Cerberus (GPU training, 192.168.1.242) --- 2x RTX 5090, 64 GB VRAM &
        mock-cerberus container &
        k3d-hydra-dev-agent-1 (label: \texttt{gpu-role=training}) \\
    \bottomrule
\end{tabularx}
\end{center}

\subsection{Service Mapping}

\begin{center}
\small
\begin{tabularx}{\textwidth}{l X X}
    \toprule
    \textbf{Production} & \textbf{Dev (Docker Compose)} & \textbf{Dev (k3d)} \\
    \midrule
    Traefik v3.6 &
        Traefik container (Docker provider) &
        Traefik deployment (K8s CRD provider) \\
    \addlinespace
    Azure AD SAML SSO &
        mock-saml-idp (SimpleSAMLphp) &
        mock-saml-idp pod \\
    \addlinespace
    SQLite (hydra.db) &
        SQLite in container volume &
        SQLite in PersistentVolumeClaim \\
    \addlinespace
    RKE2 Kubernetes &
        Docker Compose networking &
        k3d (k3s in Docker) \\
    \bottomrule
\end{tabularx}
\end{center}

\subsection{Network Topology}

In Docker Compose mode, all containers share a single Docker bridge
network. Traefik listens on the host's ports 80 and 443 and routes traffic
based on the \texttt{Host} header.

In k3d mode, the k3d cluster creates its own Docker network. Port mapping
is configured in \cmdref{k8s/dev/k3d-config.yaml} to forward host ports
80 and 443 into the k3d load balancer, which in turn routes traffic to the
Traefik ingress controller inside the cluster.

%%%%%%%%%%%%%%%%%%%%%%%%%%%%%%%%%%%%%%%%%%%%%%%%%%%%%%%%%%%%%%%%%%%%%%%%%%%
%% 4. KUBERNETES DEVELOPMENT (k3d)
%%%%%%%%%%%%%%%%%%%%%%%%%%%%%%%%%%%%%%%%%%%%%%%%%%%%%%%%%%%%%%%%%%%%%%%%%%%
\section{Kubernetes Development (k3d)}
\label{sec:k3d}

For a more production-faithful development experience, use the k3d mode.
This creates a local Kubernetes cluster with three nodes (one server and
two agents) that mirrors the production topology.

\subsection{Automated Setup with Make}

\begin{lstlisting}[style=bash]
# Step 1: Create the k3d cluster
cd hydra-saml-auth/dev
make k8s-setup

# Step 2: Verify nodes match production topology
kubectl get nodes --show-labels

# Step 3: Deploy all services
make k8s-start

# Step 4: Check overall status
make k8s-status

# Step 5: View logs from all pods
make k8s-logs
\end{lstlisting}

\notebox{The \cmdref{make k8s-setup} target creates a k3d cluster named
\texttt{hydra-dev} with node labels that mirror the production cluster.
It also applies Traefik CRDs and base namespace manifests.}

\subsection{Manual kubectl Deployment}

If you prefer to run each step manually (or need to debug a specific
stage), follow these commands:

\subsubsection{Create the k3d Cluster}

\begin{lstlisting}[style=bash]
# Create the 3-node cluster from the config file
k3d cluster create hydra-dev \
    --config k8s/dev/k3d-config.yaml
\end{lstlisting}

Verify the nodes are ready:

\begin{lstlisting}[style=bash]
kubectl get nodes -o wide
# Expected output:
#   k3d-hydra-dev-server-0   Ready   control-plane,master
#   k3d-hydra-dev-agent-0    Ready   <none>
#   k3d-hydra-dev-agent-1    Ready   <none>
\end{lstlisting}

\subsubsection{Apply Traefik CRDs}

\begin{lstlisting}[style=bash]
kubectl apply -f \
    https://raw.githubusercontent.com/traefik/traefik/v2.11/docs/content/reference/dynamic-configuration/kubernetes-crd-definition-v1.yml
\end{lstlisting}

\subsubsection{Apply Base Manifests}

\begin{lstlisting}[style=bash]
# Namespaces
kubectl apply -f k8s/base/namespace.yaml

# RBAC roles and bindings
kubectl apply -f k8s/base/rbac/

# Storage classes and PVCs
kubectl apply -f k8s/base/storage/
\end{lstlisting}

\subsubsection{Deploy Traefik}

\begin{lstlisting}[style=bash]
kubectl apply -f k8s/components/traefik/deployment.yaml
kubectl apply -f k8s/components/traefik/service.yaml
\end{lstlisting}

\subsubsection{Deploy hydra-auth}

\begin{lstlisting}[style=bash]
kubectl apply -f k8s/components/hydra-auth/
\end{lstlisting}

\subsubsection{Verify All Pods}

\begin{lstlisting}[style=bash]
kubectl get pods -n hydra-system
# All pods should show Running or Completed status

kubectl get pods -n hydra-students
# Should be empty initially (student containers are
# created on demand)
\end{lstlisting}

\subsection{Stopping and Cleaning Up}

\begin{lstlisting}[style=bash]
# Stop the k3d cluster (preserves state)
make k8s-stop

# Delete the k3d cluster entirely
make k8s-clean
# Or manually:
k3d cluster delete hydra-dev
\end{lstlisting}

%%%%%%%%%%%%%%%%%%%%%%%%%%%%%%%%%%%%%%%%%%%%%%%%%%%%%%%%%%%%%%%%%%%%%%%%%%%
%% 5. MAKEFILE REFERENCE
%%%%%%%%%%%%%%%%%%%%%%%%%%%%%%%%%%%%%%%%%%%%%%%%%%%%%%%%%%%%%%%%%%%%%%%%%%%
\section{Makefile Reference}
\label{sec:makefile}

All targets are invoked from the \cmdref{hydra-saml-auth/dev/} directory.

\subsection{Docker Compose Targets}

\begin{center}
\begin{tabularx}{\textwidth}{l X}
    \toprule
    \textbf{Target} & \textbf{Description} \\
    \midrule
    \texttt{make setup} &
        One-time setup: generates JWT key pair, creates \texttt{.env} from
        \texttt{.env.example}, adds hosts file entries for
        \texttt{hydra.local} and \texttt{gpt.hydra.local}. \\
    \addlinespace
    \texttt{make start} &
        Starts all Docker Compose services in detached mode using
        \texttt{docker-compose.dev.yml}. \\
    \addlinespace
    \texttt{make stop} &
        Stops all running containers without removing volumes. \\
    \addlinespace
    \texttt{make restart} &
        Equivalent to \texttt{make stop \&\& make start}. \\
    \addlinespace
    \texttt{make logs} &
        Tails logs from all running containers. \\
    \addlinespace
    \texttt{make clean} &
        Stops all containers and removes volumes (full reset). \\
    \addlinespace
    \texttt{make status} &
        Shows running containers and their health status. \\
    \addlinespace
    \texttt{make build} &
        Rebuilds all custom images without cache. \\
    \bottomrule
\end{tabularx}
\end{center}

\subsection{Kubernetes (k3d) Targets}

\begin{center}
\begin{tabularx}{\textwidth}{l X}
    \toprule
    \textbf{Target} & \textbf{Description} \\
    \midrule
    \texttt{make k8s-setup} &
        Creates the k3d cluster from \texttt{k8s/dev/k3d-config.yaml},
        applies CRDs, namespaces, RBAC, and storage manifests. \\
    \addlinespace
    \texttt{make k8s-start} &
        Deploys all Kubernetes workloads (Traefik, hydra-auth, mock-saml-idp,
        mock metrics agents). \\
    \addlinespace
    \texttt{make k8s-stop} &
        Stops the k3d cluster (preserves state for resume). \\
    \addlinespace
    \texttt{make k8s-status} &
        Shows nodes, pods, services, and ingresses across all namespaces. \\
    \addlinespace
    \texttt{make k8s-logs} &
        Tails logs from all pods in \texttt{hydra-system} namespace. \\
    \addlinespace
    \texttt{make k8s-clean} &
        Deletes the entire k3d cluster (full reset). \\
    \bottomrule
\end{tabularx}
\end{center}

%%%%%%%%%%%%%%%%%%%%%%%%%%%%%%%%%%%%%%%%%%%%%%%%%%%%%%%%%%%%%%%%%%%%%%%%%%%
%% 6. TROUBLESHOOTING
%%%%%%%%%%%%%%%%%%%%%%%%%%%%%%%%%%%%%%%%%%%%%%%%%%%%%%%%%%%%%%%%%%%%%%%%%%%
\section{Troubleshooting}
\label{sec:troubleshooting}

\subsection{Hosts File Not Writable}

\textbf{Symptom:} \cmdref{make setup} fails with a permission error when
writing to the hosts file.

\textbf{Fix:} Manually add the following entries:

\begin{lstlisting}[style=bash]
# On Linux/macOS: /etc/hosts
# On Windows: C:\Windows\System32\drivers\etc\hosts
# (edit as Administrator)

127.0.0.1   hydra.local
127.0.0.1   gpt.hydra.local
\end{lstlisting}

On Linux/macOS, you can also run:

\begin{lstlisting}[style=bash]
echo "127.0.0.1 hydra.local gpt.hydra.local" \
    | sudo tee -a /etc/hosts
\end{lstlisting}

\subsection{Port Conflicts (80, 443, 6969)}

\textbf{Symptom:} Containers fail to start with ``port already in use''
errors.

\textbf{Fix:} Identify and stop the conflicting process:

\begin{lstlisting}[style=bash]
# Linux/macOS
sudo lsof -i :80
sudo lsof -i :443
sudo lsof -i :6969

# Kill the conflicting process
sudo kill -9 <PID>
\end{lstlisting}

Common culprits: Apache (\texttt{apache2}), Nginx (\texttt{nginx}), or
another instance of Traefik. On macOS, AirPlay Receiver uses port 5000;
on Windows, IIS may bind to port 80.

\subsection{Docker Socket Permissions on Linux}

\textbf{Symptom:} \cmdref{docker: permission denied} when running Docker
commands without \texttt{sudo}.

\textbf{Fix:}

\begin{lstlisting}[style=bash]
sudo usermod -aG docker $USER
newgrp docker
# Or log out and log back in
\end{lstlisting}

\subsection{WSL2 Memory Limits}

\textbf{Symptom:} Containers are killed or become unresponsive on Windows
with WSL2.

\textbf{Fix:} WSL2 defaults to using up to 50\% of host RAM but may not be
enough for the full stack. Edit \cmdref{\%USERPROFILE\%\textbackslash.wslconfig}:

\begin{lstlisting}[style=yaml]
[wsl2]
memory=8GB
processors=4
swap=4GB
\end{lstlisting}

Then restart WSL2:

\begin{lstlisting}[style=bash]
wsl --shutdown
\end{lstlisting}

\subsection{k3d Cluster Won't Start}

\textbf{Symptom:} \cmdref{k3d cluster create} hangs or fails.

\textbf{Fix:}

\begin{enumerate}[leftmargin=2em]
    \item Ensure Docker is running: \cmdref{docker info}
    \item Check for existing clusters: \cmdref{k3d cluster list}
    \item Delete stale clusters: \cmdref{k3d cluster delete hydra-dev}
    \item Check Docker resource limits (at least 4 GB RAM, 2 CPUs).
    \item Retry: \cmdref{make k8s-setup}
\end{enumerate}

\subsection{SAML Callback URL Mismatch}

\textbf{Symptom:} After authenticating with the mock IdP, you are
redirected to a wrong URL or get a SAML response error.

\textbf{Fix:}

\begin{enumerate}[leftmargin=2em]
    \item Ensure your hosts file has the correct entries
          (\cmdref{hydra.local} pointing to \cmdref{127.0.0.1}).
    \item Verify the SAML callback URL in \cmdref{dev/.env}:
\begin{lstlisting}[style=bash]
SAML_CALLBACK_URL=http://hydra.local/auth/saml/callback
\end{lstlisting}
    \item Ensure the mock IdP configuration in
          \cmdref{dev/saml/} matches the callback URL.
    \item Restart services: \cmdref{make restart}
\end{enumerate}

\subsection{Student Container Creation Fails in Dev}

\textbf{Symptom:} Clicking ``Create Container'' on the dashboard returns
an error or the container never appears.

\textbf{Fix:}

\begin{itemize}[leftmargin=2em]
    \item \textbf{Docker Compose mode:} Container creation goes through
          Docker socket. Verify the hydra-saml-auth container has the
          Docker socket mounted:
\begin{lstlisting}[style=bash]
docker inspect hydra-saml-auth | grep -A2 "docker.sock"
\end{lstlisting}
    \item \textbf{k3d mode:} Container creation uses kubectl against the
          k3d cluster. Ensure the \texttt{hydra-system} service account
          has the correct RBAC permissions:
\begin{lstlisting}[style=bash]
kubectl auth can-i create pods -n hydra-students \
    --as=system:serviceaccount:hydra-system:hydra-auth
\end{lstlisting}
    \item Check logs: \cmdref{make logs} or \cmdref{make k8s-logs}
\end{itemize}

%%%%%%%%%%%%%%%%%%%%%%%%%%%%%%%%%%%%%%%%%%%%%%%%%%%%%%%%%%%%%%%%%%%%%%%%%%%
%% 7. DIFFERENCES FROM PRODUCTION
%%%%%%%%%%%%%%%%%%%%%%%%%%%%%%%%%%%%%%%%%%%%%%%%%%%%%%%%%%%%%%%%%%%%%%%%%%%
\section{Differences from Production}
\label{sec:differences}

The development environment intentionally simplifies several aspects of
the production cluster. Understanding these differences is important to
avoid confusion and to know what cannot be tested locally.

\subsection{No Real GPU Access}

The production cluster has dedicated GPU nodes:
\begin{itemize}[leftmargin=2em]
    \item \textbf{Chimera:} 3x NVIDIA RTX 3090 (72 GB VRAM total) for
          inference via Ollama and OpenWebUI.
    \item \textbf{Cerberus:} 2x NVIDIA RTX 5090 (64 GB VRAM total) for
          training workloads.
\end{itemize}

In development, GPU metrics are provided by mock agents
(\texttt{mock-chimera} and \texttt{mock-cerberus}) that return static
JSON payloads simulating GPU utilization, temperature, and memory usage.
Ollama runs on CPU in dev mode, so inference will be significantly slower.

\subsection{No ZFS RAID-10 Storage}

Production uses a 21 TB ZFS RAID-10 array on Hydra for persistent storage.
The development environment uses:

\begin{itemize}[leftmargin=2em]
    \item \textbf{Docker Compose:} Docker volumes (local filesystem).
    \item \textbf{k3d:} \texttt{local-path-provisioner} for
          PersistentVolumeClaims (data stored inside Docker volumes).
\end{itemize}

\subsection{No Real SAML SSO}

Production authenticates against Azure AD via SAML 2.0 (SUNY SSO). The
development environment uses a mock SAML IdP (SimpleSAMLphp) with
preconfigured test users. The SAML flow is identical, but the IdP is local.

\subsection{No SSH / SSHPiper}

In production, students connect to their containers via SSH through an
SSHPiper proxy. In development:

\begin{itemize}[leftmargin=2em]
    \item \textbf{Docker Compose:} Access containers via
          \cmdref{docker exec -it <container> /bin/bash}
    \item \textbf{k3d:} Access pods via
          \cmdref{kubectl exec -it <pod> -n hydra-students -- /bin/bash}
\end{itemize}

\subsection{Traefik Provider Differences}

\begin{itemize}[leftmargin=2em]
    \item \textbf{Docker Compose:} Traefik uses the \textbf{Docker provider}
          and reads labels from container definitions.
    \item \textbf{k3d:} Traefik uses the \textbf{Kubernetes CRD provider}
          and reads IngressRoute resources.
    \item \textbf{Production:} Traefik v3.6 uses the Kubernetes CRD
          provider on the RKE2 cluster.
\end{itemize}

\subsection{Summary Table}

\begin{center}
\small
\begin{tabularx}{\textwidth}{l l X}
    \toprule
    \textbf{Feature} & \textbf{Production} & \textbf{Development} \\
    \midrule
    GPU Access & Real NVIDIA GPUs &
        Mocked via metrics agents; Ollama on CPU \\
    \addlinespace
    Storage & ZFS RAID-10 (21 TB) &
        Docker volumes / local-path-provisioner \\
    \addlinespace
    Authentication & Azure AD SAML 2.0 &
        Mock SimpleSAMLphp IdP \\
    \addlinespace
    SSH Access & SSHPiper proxy &
        docker exec / kubectl exec \\
    \addlinespace
    Kubernetes & RKE2 (3 physical nodes) &
        k3d (3 Docker containers) \\
    \addlinespace
    Traefik & v3.6, K8s CRD provider &
        Docker provider (Compose) / K8s CRD (k3d) \\
    \addlinespace
    Monitoring & Prometheus + Grafana &
        Not included (optional add-on) \\
    \bottomrule
\end{tabularx}
\end{center}

%%%%%%%%%%%%%%%%%%%%%%%%%%%%%%%%%%%%%%%%%%%%%%%%%%%%%%%%%%%%%%%%%%%%%%%%%%%
%% APPENDIX
%%%%%%%%%%%%%%%%%%%%%%%%%%%%%%%%%%%%%%%%%%%%%%%%%%%%%%%%%%%%%%%%%%%%%%%%%%%
\newpage
\appendix

\section{Hosts File Entries}
\label{app:hosts}

The following entries must be present in your hosts file for the dev
environment to work correctly:

\begin{lstlisting}
127.0.0.1   hydra.local
127.0.0.1   gpt.hydra.local
\end{lstlisting}

\textbf{File locations:}
\begin{itemize}[leftmargin=2em]
    \item \textbf{Linux / macOS:} \cmdref{/etc/hosts}
    \item \textbf{Windows:}
          \cmdref{C:\textbackslash Windows\textbackslash System32\textbackslash
          drivers\textbackslash etc\textbackslash hosts}
\end{itemize}

\section{Environment Variables}
\label{app:envvars}

The \cmdref{dev/.env} file is generated by \cmdref{make setup}. Key
variables include:

\begin{center}
\begin{tabularx}{\textwidth}{l X}
    \toprule
    \textbf{Variable} & \textbf{Description} \\
    \midrule
    \texttt{NODE\_ENV} & Set to \texttt{development} \\
    \texttt{PORT} & Application port (default: \texttt{6969}) \\
    \texttt{JWT\_SECRET} & Auto-generated JWT signing secret \\
    \texttt{SAML\_ENTRY\_POINT} & Mock IdP SSO URL \\
    \texttt{SAML\_CALLBACK\_URL} & SAML assertion consumer service URL \\
    \texttt{SAML\_ISSUER} & SAML service provider entity ID \\
    \texttt{DB\_PATH} & Path to SQLite database file \\
    \texttt{RUNTIME\_MODE} & \texttt{docker} or \texttt{kubernetes} \\
    \bottomrule
\end{tabularx}
\end{center}

\section{Useful Commands}
\label{app:commands}

\begin{lstlisting}[style=bash]
# View all running containers (Docker Compose)
docker compose -f dev/docker-compose.dev.yml ps

# Rebuild a single service
docker compose -f dev/docker-compose.dev.yml build hydra-auth

# Shell into the auth container
docker exec -it hydra-saml-auth /bin/sh

# View k3d cluster info
k3d cluster list
kubectl cluster-info

# Port-forward a specific service (k3d)
kubectl port-forward svc/hydra-auth 6969:6969 -n hydra-system

# Check resource usage in k3d
kubectl top nodes
kubectl top pods -n hydra-system
\end{lstlisting}

\end{document}
